\documentclass[12pt]{article}
\usepackage[utf8]{inputenc}
\usepackage[T1]{fontenc}
\usepackage[margin=1in]{geometry}
\usepackage{setspace}
\usepackage{xcolor}
\usepackage[colorlinks=true, urlcolor=blue!60!black, linkcolor=blue!60!black]{hyperref}
\setlength{\parindent}{0em}
\setlength{\parskip}{0.6em}
\renewcommand{\baselinestretch}{1.0}
\begin{document}
\begin{center}
{\LARGE \textit{Florida Prepaid Postsecondary Education Expense Board v.\ College Savings Bank}, 527 U.S.\ 627 (1999)\par}
\vspace{0.5em}
{\large Case Brief --- Two-Page Edition\par}
\end{center}

\section{Memory Jogger}
The Patent Remedy Act's abrogation of state immunity failed \textit{City of Boerne}'s ``congruence and proportionality'' test (\href{https://www.courtlistener.com/opinion/118316/florida-prepaid-postsecondary-education-expense-board-v-college-savings/}{CourtListener}).

\section{Detailed Case Facts}
College Savings Bank patented its method for tuition-guarantee CDs and alleged that Florida Prepaid, an arm of the State, willfully infringed it. Congress's 1992 Patent Remedy Act declared that States ``shall not be immune'' from federal infringement suits.

\section{Procedural History}
After \textit{Seminole Tribe}, Florida Prepaid moved to dismiss on immunity grounds. The district court denied the motion, and the Federal Circuit affirmed, upholding the Act under \S~5.

\section{Judicial Votes}
Reversed, 5--4. Rehnquist wrote for O'Connor, Scalia, Kennedy, and Thomas; Stevens dissented, joined by Souter, Ginsburg, and Breyer.

\section{Holding}
The Act is not valid \S~5 legislation. Congress must identify a pattern of Fourteenth Amendment violations, here intentional infringement without adequate state remedies, and tailor a congruent and proportional response.

\section{Analysis of Opinions}
\textbf{Majority:} The record showed few state infringements (eight suits in 110 years). Infringement violates due process only if the State offers no adequate remedy, and negligence is no deprivation, yet Congress barely studied state remedies and swept in innocent infringement by all 50 States. Its real aim, uniformity, is an Article~I concern. \textbf{Dissent:} exclusive federal jurisdiction made state remedies inadequate; the Act changed no substantive law and scaled with actual infringement; and it was unfair to demand a record for a test announced after 1992.

\section{Examples: Future Applications}
\textbf{Same-Side:} copyright abrogation on a thin record fails (\textit{Allen v.\ Cooper}). \textbf{Opposite-Side:} the FMLA family-care provision, resting on a record of gender discrimination, is valid (\textit{Hibbs}). \textbf{Fence-Sitter:} conditioning research grants on waiver of immunity under the Spending Clause.

\section{Critique}
Textualists note that the Eleventh Amendment's text bars only diversity suits; progressive scholars see a ``juricentric Constitution.'' States enforce their own patents yet escape liability for infringing others'.

\section{Key Quotations}
\begin{quote}
``We thus held that for Congress to invoke \S~5, it must identify conduct transgressing the Fourteenth Amendment's substantive provisions, and must tailor its legislative scheme to remedying or preventing such conduct.''
\end{quote}
\noindent --- Chief Justice Rehnquist, for the Court, 527 U.S.\ at 639.

\end{document}
